\scanpage{16}
% 1–7
Soit $p:X\to Y$ un morphisme submersif de schémas formels (lisses), \unsure{supposons} $X$ muni d'une $A_*$-structure. On dit que $Y$ admet une $A_*$-structure quotient de celle de $X$ (moyennant $p$) s'il admet une \inserted{$A_*$} structure telle que $p$ devienne un $A_*$-morphisme. Celle-ci est encore unique : en effet, il existe (sans condition sur $Y$)

% 8–14
un \deleted{\unsure{sous $A$|sous-$A_*$}} sous-schéma $A_*$-formel $Y'$ de $X$, tel que le morphisme formel induit $Y'\to Y$ soit un isomorphisme. \inserted{Considérons} le composé
\diagram{\begin{tikzcd}[ampersand replacement=\&]
\& Y\arrow[dr,"\sim"] \& \\
X\arrow[ur,"p"]\arrow[rr] \& \& Y'
\end{tikzcd}}
\deleted{\unsure{alors}}, on voit alors grâce à c) que la $A_*$-structure quotient existe ssi $X'\to Y'$ est un $A_*$-morphisme, et la $A_*$-structure sur $Y$ est celle déduite de celle de $Y'$.

% 15–19
\deleted{\unsure{Une co|Une ca}} Soit $X\xrightarrow{p}Y$ un morphisme formel submersif. Alors $R=X\mathbin{\widehat{\times}}_Y X\subset X\mathbin{\widehat{\times}}X$ est un \inserted{sous-schéma formel lisse, et c'est} relation d'équivalence [stable par symétrie, contient diagonale, transitif\ldots]. \deleted{\unsure{Inversement}}\inserted{Plus précisément},

% 20–25
$X$ étant donné, il y a corr. biunivoque entre ces quotients formels $Y$, et les $R$ rel. d'équiv. formelles. Si $\deleted{A}X$ est muni d'une $A_*$-structure, alors si $Y$ admet une $A_*$-structure quotient, on voit que $R$ est un sous-schéma $A_*$-formel de $X\mathbin{\widehat{\times}}X$ i.e. il a une $A_*$-structure induite.

% 26–28
f) (Condition intégrale) La réciproque est vraie, i.e. si $R$ est une relation d'équival. formelle dans $X$, qui soit une $A_*$-sous-structure de $X\mathbin{\widehat{\times}}X$, alors
% archive sheet number: 15
